% For Section IV (conditions, parameters and setup). Moved from draft 1 of Section V on 2 Oct 2026
% to keep the results within 2.5 pages. Section V refers to it as \ref{sec:stats}.

\subsection{Measurement and statistics}\label{sec:stats}

Service is probe success $S$: the share of requests from the three probe hosts (h1, h6 and h11, which no arm may block) answered within a 0.5\,s step. A step in which the probes completed no request scores 0. ``After the event'' means steps 50 to 99 and ``before'' means steps 2 to 49, which leaves out the two reset steps. We take the worst window as the lowest $S$ over steps 52 to 99. ISASF's version starts at step 2, but in both conditions the event restarts the probe clients and leaves step 50 empty, which pins it at 0 for every arm. For recovery we use ISASF's rule: the five-step rolling mean of $S$ must return to 95\% of its mean over the 20 steps before the clock starts and stay there for five steps. We start the clock at step 51 because ISASF's regime label changes one row after the same dip.

The unit of analysis is the seed. An arm's value for a seed is the mean over that seed's five episodes, which gives ten values per arm and condition. (The static rule has no weights, so its ``seed'' is the block of five environment seeds it shares with the learned arms, and all four arms meet the same attackers, surge hosts and benign rates.) We report the interquartile mean (IQM) of the ten values, the mean of the middle half once the lowest and highest quarters are dropped, with a percentile bootstrap 95\% interval over the seeds (as in~\cite{agarwal2021}, 10{,}000 resamples). A2 is compared with each other arm by an exact paired Wilcoxon signed-rank test over the seed pairs~\cite{wilcoxon1945}, which asks whether one arm sits above the other consistently from seed to seed; seeds whose two values are equal drop out. Holm's method~\cite{holm1979} adjusts the three $p$-values of each parameter, and a difference counts as significant when the adjusted value $p_H$ is below 0.05. With ten seeds the smallest attainable $p$ is $2/2^{10}=0.00195$, which Holm multiplies by three, so $p_H=0.0059$ means that A2 fell on the same side of the other arm in all ten seeds.
